\RequirePackage{etex}
\documentclass[runningheads,envcountsame,a4paper]{llncs}

\usepackage[T1]{fontenc} % fix missing font cmtt
\usepackage{amsmath}
\usepackage{amssymb} % Vdash
\usepackage{graphicx} % rotatebox
\usepackage{stmaryrd} % llparenthesis
\usepackage{anyfontsize} % workaround for font size difference warning
\usepackage{listings}
\usepackage{tikz}
\usetikzlibrary{calc,fit,tikzmark,plotmarks,arrows.meta,positioning,overlay-beamer-styles}
\usepackage[caption=false]{subfig}

\usepackage{cancel} % slash over symbol
\usepackage{hyperref}
\renewcommand\UrlFont{\color{blue}\rmfamily}

\usepackage{xcolor}
\definecolor{hazelgreen}{RGB}{7,63,36}
\definecolor{hazellightgreen}{RGB}{103,138,97}
\definecolor{hazelyellow}{RGB}{245,222,179}
\definecolor{hazellightyellow}{RGB}{254,254,234}

\newcommand{\highlight}[1]{\colorbox{yellow}{$\displaystyle #1$}}

\newcommand{\todo}[1]{{\color{red} TODO: #1}}

%% Joshua Dunfield macros
\def\OPTIONConf{1}
\usepackage{jdunfield}
\usepackage{rulelinks} % hyperlink of rule name
% !TEX root = ./patterns-paper.tex

% reverse Vdash
\newcommand{\dashV}{\mathbin{\rotatebox[origin=c]{180}{$\Vdash$}}}

% Violet hotdogs; highlight color helps distinguish them
\newcommand{\llparenthesiscolor}{\textcolor{violet}{\llparenthesis}}
\newcommand{\rrparenthesiscolor}{\textcolor{violet}{\rrparenthesis}}

% HTyp and HExp
\newcommand{\hcomplete}[1]{#1~\mathsf{complete}}

% HTyp
\newcommand{\htau}{\dot{\tau}}
\newcommand{\tarr}[2]{\inparens{#1 \rightarrow #2}}
\newcommand{\tarrnp}[2]{#1 \rightarrow #2}
\newcommand{\trul}[2]{\inparens{#1 \Rightarrow #2}}
\newcommand{\tnum}{\mathtt{num}}
\newcommand{\tehole}{\llparenthesiscolor\rrparenthesiscolor}
\newcommand{\tsum}[2]{\inparens{{#1} + {#2}}}
\newcommand{\tprod}[2]{\inparens{{#1} \times {#2}}}
\newcommand{\tunit}{\mathtt{1}}
\newcommand{\tvoid}{\mathtt{0}}

\newcommand{\tcompat}[2]{#1 \sim #2}
\newcommand{\tincompat}[2]{#1 \nsim #2}

% HExp
\newcommand{\hexp}{\dot{e}}
\newcommand{\hlam}[3]{\inparens{\lambda #1:#2.#3}}
\newcommand{\hap}[2]{#1(#2)}
\newcommand{\hapP}[2]{(#1)~(#2)} % Extra paren around function term
\newcommand{\hnum}[1]{\underline{#1}}
\newcommand{\hadd}[2]{\inparens{#1 + #2}}
\newcommand{\hpair}[2]{\inparens{#1 , #2}}
\newcommand{\htriv}{()}
\newcommand{\hehole}{\llparenthesiscolor\rrparenthesiscolor}
\newcommand{\hhole}[1]{\llparenthesiscolor#1\rrparenthesiscolor}
\newcommand{\hindet}[1]{\lceil#1\rceil}
\newcommand{\hinj}[2]{\mathtt{inj}_{#1}({#2})}
\newcommand{\hinl}[2]{\mathtt{inl}_{#1}({#2})}
\newcommand{\hinr}[2]{\mathtt{inr}_{#1}({#2})}
\newcommand{\hinlp}[1]{\mathtt{inl}(#1)}
\newcommand{\hinrp}[1]{\mathtt{inr}(#1)}
\newcommand{\hmatch}[2]{\mathtt{match}(#1) \{#2\}}
\newcommand{\hcase}[5]{\mathtt{case}({#1},{#2}.{#3},{#4}.{#5})}
\newcommand{\hrules}[2]{\inparens{#1 \mid #2}}
\newcommand{\hrul}[2]{#1 \Rightarrow #2}

\newcommand{\hGamma}{\dot{\Gamma}}
\newcommand{\domof}[1]{\text{dom}(#1)}
\newcommand{\hsyn}[3]{#1 \vdash #2 \Rightarrow #3}
\newcommand{\hana}[3]{#1 \vdash #2 \Leftarrow #3}
\newcommand{\hexptyp}[3]{#1 \vdash #2 : #3}
\newcommand{\hpattyp}[3]{#1 : #2 \dashV #3}
\newcommand{\hpatmatch}[3]{#1 \vartriangleright #2 \dashV #3}
\newcommand{\hval}[1]{#1 ~\mathtt{val}}
\newcommand{\herr}[1]{#1 ~\mathtt{err}}

% ZTyp and ZExp
\newcommand{\zlsel}[1]{{\bowtie}{#1}}
\newcommand{\zrsel}[1]{{#1}{\bowtie}}
\newcommand{\zwsel}[1]{
  \setlength{\fboxsep}{0pt}
  \colorbox{green!10!white!100}{
    \ensuremath{{{\textcolor{Green}{{\hspace{-2px}\triangleright}}}}{#1}{\textcolor{Green}{\triangleleft{\vphantom{\tehole}}}}}}
}

\newcommand{\removeSel}[1]{#1^{\diamond}}

% ZTyp
\newcommand{\ztau}{\hat{\tau}}

% ZExp
\newcommand{\zexp}{\hat{e}}

% Direction
\newcommand{\dParent}{\mathtt{parent}}
\newcommand{\dChildn}[1]{\mathtt{child}~\mathtt{{#1}}}
\newcommand{\dChildnm}[1]{\mathtt{child}~{#1}}

% Action
\newcommand{\aMove}[1]{\mathtt{move}~#1}
	\newcommand{\zrightmost}[1]{\mathsf{rightmost}(#1)}
	\newcommand{\zleftmost}[1]{\mathsf{leftmost}(#1)}
\newcommand{\aSelect}[1]{\mathtt{sel}~#1}
\newcommand{\aDel}{\mathtt{del}}
\newcommand{\aReplace}[1]{\mathtt{replace}~#1}
\newcommand{\aConstruct}[1]{\mathtt{construct}~#1}
\newcommand{\aConstructx}[1]{#1}
\newcommand{\aFinish}{\mathtt{finish}}

\newcommand{\performAna}[5]{#1 \vdash #2 \xlongrightarrow{#4} #5 \Leftarrow #3}
\newcommand{\performAnaI}[5]{#1 \vdash #2 \xlongrightarrow{#4}\hspace{-3px}{}^{*}~ #5 \Leftarrow #3}
\newcommand{\performSyn}[6]{#1 \vdash #2 \Rightarrow #3 \xlongrightarrow{#4} #5 \Rightarrow #6}
\newcommand{\performSynI}[6]{#1 \vdash #2 \Rightarrow #3 \xlongrightarrow{#4}\hspace{-3px}{}^{*}~ #5 \Rightarrow #6}
\newcommand{\performTyp}[3]{#1 \xlongrightarrow{#2} #3}
\newcommand{\performTypI}[3]{#1 \xlongrightarrow{#2}\hspace{-3px}{}^{*}~#3}

\newcommand{\performMove}[3]{#1 \xlongrightarrow{#2} #3}
\newcommand{\performDel}[2]{#1 \xlongrightarrow{\aDel} #2}

% Form
\newcommand{\farr}{\mathtt{arrow}}
\newcommand{\fnum}{\mathtt{num}}
\newcommand{\fsum}{\mathtt{sum}}

\newcommand{\fasc}{\mathtt{asc}}
\newcommand{\fvar}[1]{\mathtt{var}~#1}
\newcommand{\flam}[1]{\mathtt{lam}~#1}
\newcommand{\fap}{\mathtt{ap}}
% \newcommand{\farg}{\mathtt{arg}}
\newcommand{\fnumlit}[1]{\mathtt{lit}~#1}
\newcommand{\fplus}{\mathtt{plus}}
\newcommand{\fhole}{\mathtt{hole}}
\newcommand{\fnehole}{\mathtt{nehole}}

\newcommand{\finj}[1]{\mathtt{inj}~#1}
\newcommand{\fcase}[2]{\mathtt{case}~#1~#2}

% Talk about formal rules in example
\newcommand{\refrule}[1]{\textrm{Rule~(#1)}}

\newcommand{\herase}[1]{\left|#1\right|_\textsf{erase}}

\newcommand{\arrmatch}[2]{#1 \blacktriangleright_{\rightarrow} #2}


\newcommand{\TABperformAna}[5]{#1 \vdash & #2                & \xlongrightarrow{#4} & #5 & \Leftarrow #3}
\newcommand{\TABperformSyn}[6]{#1 \vdash & #2 \Rightarrow #3 & \xlongrightarrow{#4} & #5 \Rightarrow #6}
\newcommand{\TABperformTyp}[3]{& #1 & \xlongrightarrow{#2} & #3}

\newcommand{\TABperformMove}[3]{#1 & \xlongrightarrow{#2} & #3}
\newcommand{\TABperformDel}[2]{#1 \xlongrightarrow{\aDel} #2}

\newcommand{\sumhasmatched}[2]{#1 \mathrel{\textcolor{black}{\blacktriangleright_{+}}} #2}

\newcommand{\subminsyn}[1]{\mathsf{submin}_{\Rightarrow}(#1)}
\newcommand{\subminana}[1]{\mathsf{submin}_{\Leftarrow}(#1)}


\newcommand{\inparens}[1]{{\color{gray}(}#1{\color{gray})}}

%% rule names for appendix
\newcommand{\rname}[1]{\textsc{#1}}
\newcommand{\gap}{\vspace{7pt}}


% !TEX root = pattern-paper.tex
%% Statics

% typing of internal expressions
\newrulecommand{TVar}{TVar}
\newrulecommand{TNum}{TNum}
\newrulecommand{TLam}{TLam}
\newrulecommand{TAp}{TAp}
\newrulecommand{TPair}{TPair}
\newrulecommand{TPrl}{TPrl}
\newrulecommand{TPrr}{TPrr}
\newrulecommand{TInl}{TInl}
\newrulecommand{TInr}{TInr}
\newrulecommand{TMatchZPre}{TMatchZPre}
\newrulecommand{TMatchNZPre}{TMatchNZPre}
\newrulecommand{TEHole}{TEHole}
\newrulecommand{THole}{THole}

% typing of rule and rules
\newrulecommand{TRule}{TRule}
\newrulecommand{TOneRules}{TOneRules}
\newrulecommand{TRules}{TRules}

% typing of pattern
\newrulecommand{PTVar}{PTVar}
\newrulecommand{PTNum}{PTNum}
\newrulecommand{PTWild}{PTWild}
\newrulecommand{PTPair}{PTPair}
\newrulecommand{PTInl}{PTInl}
\newrulecommand{PTInr}{PTInr}
\newrulecommand{PTEHole}{PTEHole}
\newrulecommand{PTHole}{PTHole}

%% Dynamics

% typing of substitution
\newrulecommand{STEmpty}{STEmpty}
\newrulecommand{STExtend}{STExtend}

% refutable pattern
\newrulecommand{RNum}{RNum}
\newrulecommand{REHole}{REHole}
\newrulecommand{RHole}{RHole}
\newrulecommand{RInl}{RInl}
\newrulecommand{RInr}{RInr}
\newrulecommand{RPairL}{RPairL}
\newrulecommand{RPairR}{RPairR}

% match
\newrulecommand{MVar}{MVar}
\newrulecommand{MNum}{MNum}
\newrulecommand{MWild}{MWild}
\newrulecommand{MPair}{MPair}
\newrulecommand{MEHolePair}{MEHolePair}
\newrulecommand{MHolePair}{MHolePair}
\newrulecommand{MApPair}{MApPair}
\newrulecommand{MMatchPair}{MMatchPair}
\newrulecommand{MPrlPair}{MPrlPair}
\newrulecommand{MPrrPair}{MPrrPair}
\newrulecommand{MInl}{MInl}
\newrulecommand{MInr}{MInr}

% may match
\newrulecommand{MMEHole}{MMEHole}
\newrulecommand{MMHole}{MMHole}
\newrulecommand{MMExpEHole}{MMExpEHole}
\newrulecommand{MMExpHole}{MMExpHole}
\newrulecommand{MMAp}{MMAp}
\newrulecommand{MMMatch}{MMMatch}
\newrulecommand{MMPrl}{MMPrl}
\newrulecommand{MMPrr}{MMPrr}
\newrulecommand{MMPairL}{MMPairL}
\newrulecommand{MMPairR}{MMPairR}
\newrulecommand{MMPair}{MMPair}
\newrulecommand{MMInl}{MMInl}
\newrulecommand{MMInr}{MMInr}

% not match
\newrulecommand{NMNum}{NMNum}
\newrulecommand{NMPairL}{NMPairL}
\newrulecommand{NMPairR}{NMPairR}
\newrulecommand{NMConfL}{NMConfL}
\newrulecommand{NMConfR}{NMConfR}
\newrulecommand{NMInl}{NMInl}
\newrulecommand{NMInr}{NMInr}

% value
\newrulecommand{VNum}{VNum}
\newrulecommand{VLam}{VLam}
\newrulecommand{VPair}{VPair}
\newrulecommand{VInl}{VInl}
\newrulecommand{VInr}{VInr}

% indeterminate
\newrulecommand{IEHole}{IEHole}
\newrulecommand{IHole}{IHole}
\newrulecommand{IAp}{IAp}
\newrulecommand{IPairL}{IPairL}
\newrulecommand{IPairR}{IPairR}
\newrulecommand{IPair}{IPair}
\newrulecommand{IPrl}{IPrl}
\newrulecommand{IPrr}{IPrr}
\newrulecommand{IInl}{IInl}
\newrulecommand{IInr}{IInr}
\newrulecommand{IMatch}{IMatch}

% final
\newrulecommand{FVal}{FVal}
\newrulecommand{FIndet}{FIndet}

% step
\newrulecommand{ITHole}{ITHole}
\newrulecommand{ITApFun}{ITApFun}
\newrulecommand{ITApArg}{ITApArg}
\newrulecommand{ITAp}{ITAp}
\newrulecommand{ITPairL}{ITPairL}
\newrulecommand{ITPairR}{ITPairR}
\newrulecommand{ITPrl}{ITPrl}
\newrulecommand{ITPrr}{ITPrr}
\newrulecommand{ITInl}{ITInl}
\newrulecommand{ITInr}{ITInr}
\newrulecommand{ITExpMatch}{ITExpMatch}
\newrulecommand{ITSuccMatch}{ITSuccMatch}
\newrulecommand{ITFailMatch}{ITFailMatch}

% typing of constraints
\newrulecommand{CTTruth}{CTTruth}
\newrulecommand{CTFalsity}{CTFalsity}
\newrulecommand{CTUnknown}{CTUnknown}
\newrulecommand{CTNum}{CTNum}
\newrulecommand{CTNotNum}{CTNotNum}
\newrulecommand{CTAnd}{CTAnd}
\newrulecommand{CTOr}{CTOr}
\newrulecommand{CTInl}{CTInl}
\newrulecommand{CTInr}{CTInr}
\newrulecommand{CTPair}{CTPair}

% satisfaction judgment
\newrulecommand{CSTruth}{CSTruth}
\newrulecommand{CSNum}{CSNum}
\newrulecommand{CSNotNum}{CSNotNum}
\newrulecommand{CSAnd}{CSAnd}
\newrulecommandONE{CSOrL}{CSOrL}
\newrulecommandONE{CSOrR}{CSOrR}
\newrulecommand{CSInl}{CSInl}
\newrulecommand{CSInr}{CSInr}
\newrulecommand{CSPair}{CSPair}
\newrulecommand{CSEHolePair}{CSEHolePair}
\newrulecommand{CSHolePair}{CSHolePair}
\newrulecommand{CSApPair}{CSApPair}
\newrulecommand{CSMatchPair}{CSMatchPair}
\newrulecommand{CSPrlPair}{CSPrlPair}
\newrulecommand{CSPrrPair}{CSPrrPair}

% maybe satisfaction judgment
\newrulecommand{CMSUnknown}{CMSUnknown}
\newrulecommand{CMSExpEHole}{CMSExpEHole}
\newrulecommand{CMSExpHole}{CMSExpHole}
\newrulecommand{CMSAp}{CMSAp}
\newrulecommand{CMSMatch}{CMSMatch}
\newrulecommand{CMSPrl}{CMSPrl}
\newrulecommand{CMSPrr}{CMSPrr}
\newrulecommand{CMSAndL}{CMSAndL}
\newrulecommand{CMSAndR}{CMSAndR}
\newrulecommand{CMSAnd}{CMSAnd}
\newrulecommand{CMSOrL}{CMSOrL}
\newrulecommand{CMSOrR}{CMSOrR}
\newrulecommand{CMSInl}{CMSInl}
\newrulecommand{CMSInr}{CMSInr}
\newrulecommand{CMSPairL}{CMSPairL}
\newrulecommand{CMSPairR}{CMSPairR}
\newrulecommand{CMSPair}{CMSPair}

% must or maybe satisfaction judgment
\newrulecommand{CSMSMay}{CSMSMay}
\newrulecommand{CSMSSat}{CSMSSat}

% incon
\newrulecommand{CINCTruth}{CINCTruth}
\newrulecommand{CINCFalsity}{CINCFalsity}
\newrulecommand{CINCNum}{CINCNum}
\newrulecommand{CINCNotNum}{CINCNotNum}
\newrulecommand{CINCAnd}{CINCAnd}
\newrulecommand{CINCOr}{CINCOr}
\newrulecommand{CINCInj}{CINCInj}
\newrulecommand{CINCInl}{CINCInl}
\newrulecommand{CINCInr}{CINCInr}
\newrulecommand{CINCPairL}{CINCPairL}
\newrulecommand{CINCPairR}{CINCPairR}

\begin{document}
  
\title{Pattern Matching with Typed Holes}

\author{
  Yongwei Yuan\inst{1} \and 
}
\authorrunning{Y. Yuan et al.}
\institute{
  Purdue University, West Lafayette, USA \\
  \email{yuan311@purdue.edu} \and
  University of Michigan, Ann Arbor, USA
  \email{\{dmoo,hkpotter,wyuning,comar\}@umich.edu}}


\maketitle

\begin{abstract}
  The abstract should briefly summarize the contents of the paper in
  15--250 words.
  
  \keywords{live programming  \and typed holes \and contextual modal type theory.}
\end{abstract}

\section{Introduction}
\label{sec:intro}
Most contemporary programming environments either provide meaningful feedback only for complete programs, or fall back to limited heuristic approaches when the program is incomplete.
Recent work on modeling incomplete programs as typed expressions with \emph{holes} has focused on tackling this problem. \cite{DBLP:conf/popl/OmarVHAH17} describes a static semantics for incomplete functional programs. Based on that, \cite{DBLP:journals/pacmpl/OmarVCH19} develops a dynamic semantics to evaluate such incomplete programs.
These foundational treatments are being implemented and extended in the Hazel programming environment (hazel.org).

\emph{Pattern matching} is a cornerstone of functional programming languages in the ML family. 
However, \cite{DBLP:journals/pacmpl/OmarVCH19} only supports simple case analysis on binary sum types and does not support nested patterns nor pattern holes.
This paper addresses this problem, focusing on adding full ML-style pattern matching with support for pattern holes to the Hazelnut core calculus and implementing 
this system at full scale into Hazel. Consider the examples below, which contain expression and pattern holes, denoted $\hehole{u}$ (where each hole has a unique name, $u$). 

For the match expression shown in \Eqref{eq:complete}, the scrutinee contains a hole.
By considering the match two rules in order, we observe that
no matter how the hole $u$ is filled, $\hpair{\hinl{\tnum}{\hehole{u}}}{2}$ doesn't match $\hpair{\hinrp{x}}{\_}$,
and always matches $\hpair{\_}{x}$. And thus the value of the match expression is $2$, despite the fact that the program is incomplete.
If, however, we replaced the pattern in the first rule with $(\hinlp{3}, \_)$, it would be impossible to decide whether the pattern matches
or not, and so the match expression would be \emph{indeterminate}, i.e. awaiting hole filling.
%Similarly, if we had replaced the pattern in the first rule with $(\hinlp{x}, \hehole{w})$,
%then it is impossible to tell whether the pattern hole $w$ will match $2$ or not,
%so the match would also be indeterminate.

One of the benefits of pattern matching is that it allows for static reasoning about exhaustiveness and redundancy.
However, reasoning about these matters is subtle and complicated in the presence of holes. 
For example, consider the problem of checking if the rules in \Eqref{eq:pathole-exhaustive} cover all the possibilities, \ie \emph{exhaustiveness}, or if the second rule in \Eqref{eq:pathole-redundant} can never be reached no matter how the holes are filled, \ie \emph{redundancy}. 

This abstract focuses on the formalism of full-fledged pattern matching with typed holes so that the programming environment could give this sort of feedback even when considering incomplete match expressions.

\section{Live Pattern Matching}\label{sec:examples}

\begin{figure}[ht]
  \begin{tikzpicture}[remember picture, overlay]
    \node[rotate=-90,mark size=5pt,color=hazellightgreen,shift={(-0.1,-0.2)}] at (pic cs:exp-hole-1) {\pgfuseplotmark{triangle*}};
    \node[rotate=-90,mark size=5pt,color=hazellightgreen,shift={(-0.1,-0.2)}] at (pic cs:exp-hole-2) {\pgfuseplotmark{triangle*}};
  \end{tikzpicture}

  \centering
  \hspace*{\fill}%
  \captionsetup[subfloat]{labelformat=empty}
  \subfloat[]{%
    \makebox[\width]{
      $\begin{aligned}
        &\match (1::\?) \{ \\
        &\tikzmark{exp-hole-1}|~ \hrul{\nil}{\none} \\
        &|~ \hrul{x::xs}{\some{xs}} \\
        &\}
      \end{aligned}$
    }
  }%
  $\mapsto$
  \subfloat[]{
    \makebox[1.2\width]{
      $\begin{aligned}
        &\match (1::\?) \{ \\
        &|~ \hrul{\nil}{\none} \\
        &\tikzmark{exp-hole-2}|~ \hrul{1::xs}{\some{xs}} \\
        &\}
      \end{aligned}$
    }
  }%
  $\mapsto$
  \subfloat[]{
    $\some{?}$
  }
  \hspace*{\fill}%
  \caption{Pattern Matching with Expression Holes}
  \label{fig:exp-hole-step}
\end{figure}

\begin{figure}[ht]
  \begin{tikzpicture}[remember picture, overlay]
    \node[rotate=-90,mark size=5pt,color=hazellightgreen,shift={(-0.1,-0.2)}] at (pic cs:pat-hole-1) {\pgfuseplotmark{triangle*}};
    \node[rotate=-90,mark size=5pt,color=hazellightgreen,shift={(-0.1,-0.2)}] at (pic cs:pat-hole-2) {\pgfuseplotmark{triangle*}};
  \end{tikzpicture}

  \centering
  \hspace*{\fill}%
  \captionsetup[subfloat]{labelformat=empty}
  \subfloat[]{%
    \makebox[\width]{
      $\begin{aligned}
        &\match (1::\nil) \{ \\
        &\tikzmark{pat-hole-1}|~ \hrul{\nil}{\none} \\
        &|~ \hrul{\?::xs}{\some{xs}} \\
        &\}
      \end{aligned}$
    }
  }%
  $\mapsto$
  \subfloat[]{
    \makebox[1.2\width]{
      $\begin{aligned}
        &\match (1::\nil) \{ \\
        &|~ \hrul{\nil}{\none} \\
        &\tikzmark{pat-hole-2}|~ \hrul{\?::xs}{\some{xs}} \\
        &\}
      \end{aligned}$
    }
  }%
  $\not\mapsto$(indeterminate)
  \hspace*{\fill}%
  \caption{Pattern Matching with Pattern Holes}
  \label{fig:pat-hole-step}
\end{figure}

\begin{figure}[ht]
  \centering
  \hspace*{\fill}%
  \subfloat[May or may not be exhaustive.\label{fig:may-exhaustive}]{
    \makebox[2\width]{
      $\begin{aligned}
        &\match (1::\nil) \{ \\
        &|~ \hrul{\nil}{\none} \\
        &|~ \hrul{\?::xs}{\some{xs}} \\
        &\}
      \end{aligned}$
    }
  }%
  \hfill%
  \subfloat[Must not be exhaustive.\label{fig:not-exhasutive}]{
    \makebox[2\width]{
      $\begin{aligned}
        &\match (1::\nil) \{ \\
        &|~ \hrul{\nil}{\none} \\
        &|~ \hrul{1::\?}{\?} \\
        &\}
      \end{aligned}$
    }
  }%
  \hspace*{\fill}%
  \caption{Exhaustiveness Checking in Incomplete Match Expressions}
  \label{fig:exh-hole}
\end{figure}

\begin{figure}[ht]
  \centering
  \hspace*{\fill}%
  \subfloat[May or may not be redundant.\label{fig:may-redundant}]{
    \makebox[1.3\width]{
      $\begin{aligned}
        &\match (1::2::\nil) \{ \\
        &|~ \hrul{\nil}{\none} \\
        &|~ \hrul{x::\?}{\?} \\
        &|~ \hrul{1::xs}{\some{2::xs}} \\
        &\}
      \end{aligned}$
    }
  }%
  \hfill%
  \subfloat[Must be redundant.\label{fig:redundant}]{
    \makebox[1.3\width]{
      $\begin{aligned}
        &\match (1::2::\nil) \{ \\
        &|~ \hrul{\nil}{\none} \\
        &|~ \hrul{x::xs}{\some{xs}} \\
        &|~ \hrul{\?::xs}{\some{xs}} \\
        &\}
      \end{aligned}$
    }
  }%
  \hspace*{\fill}%
  \caption{Redundancy Checking in Incomplete Match Expressions}
  \label{fig:red-hole}
\end{figure}

\section{The Core Calculus}
With the intuition developed in the previous section, we can formally specify
what we mean by pattern matching with holes and also build the type system to
give eager feedback while the user is programming.

% !TEX root= pattern-paper.tex

\begin{figure}[ht]
    \centering
    \begin{minipage}{.5\linewidth}
  $\arraycolsep=4pt\begin{array}{lll}
    e & ::= &
      x ~\vert~
      \hnum{n} \\
      & ~\vert~ &
      \hlam{x}{\tau}{e} ~\vert~
      \hap{e_1}{e_2} \\
      & ~\vert~ &
      \hpair{e_1}{e_2} ~\vert~
      \hfst{e} ~\vert~ \hsnd{e} \\
      & ~\vert~ &
      \hinl{\tau}{e} ~\vert~
      \hinr{\tau}{e} \\
      & ~\vert~ &
      \hmatch{e}{\hat{rs}} \\
      & ~\vert~ &
      \hehole{u} ~\vert~
      \hhole{e}{u} \\
    \end{array}$
    \end{minipage}%
    \begin{minipage}{.5\linewidth}
  $\arraycolsep=4pt\begin{array}{lll}
    \tau & ::= &
      \tnum ~\vert~
      \tarr{\tau_1}{\tau_2} ~\vert~
      \tprod{\tau_1}{\tau_2} ~\vert~
      \tsum{\tau_1}{\tau_2} \\
    \hat{rs} & ::= &
      \zrulsP{rs}{r}{rs} \\
    rs & ::= &
      \cdot ~\vert~ \hrulesP{r}{rs'} \\
    r & ::= &
      \hrul{p}{e} \\
    p & ::= &
      x ~\vert~
      \_ ~\vert~
      \hnum{n} ~\vert~
      \hpair{p_1}{p_2} \\
      & ~\vert~ &
      \hinlp{p} ~\vert~
      \hinrp{p} ~\vert~
      \heholep{w} ~\vert~
      \hholep{p}{w}{\tau}
    \end{array}$
    \end{minipage}
\caption{Syntax}
\label{fig:syntax}
\end{figure}

\paragraph{Overview.}
The syntax of the core calculus is given in \figurename~\ref{fig:syntax}. To
balance simplicity and power of pattern matching with holes, we use number as
the base type and include function type, product type, and sum type in the
calculus. For each type, there is a corresponding expression serving as the
introduction of the type. Besides, holes, as missing parts in a program, can be
of any type (see Sec.~\ref{sec:statics}). And the most interesting part is the
syntax of match expressions. A match expression consists of an expression,
called scrutinee, and zipper rules \todo{ref}. Zipper rules denotes a series of
rules with a pointer on the current rule being considered in runtime (see
Sec.~\ref{sec:dynamics}). As in other ML languages, each rule consists of a
pattern and an expression \todo{ref}. And just like expressions, patterns may
also involve holes and are assigned corresponding types (see Sec.~\ref{sec:statics}).

\subsection{Dynamic Semantics}\label{sec:dynamics}
Pattern matching happens during runtime and the dynamic semantics motivates the
design of the static semantics in Sec.~\ref{sec:statics}. And thus, we start by
taking a step further from the examples shown in
\listfigurename~\ref{fig:exp-hole-step},\ref{fig:pat-hole-step}, and formalize
the described heuristic.

% !TEX root= pattern-paper.tex

\begin{figure}[ht]

\judgbox{\rmpointer{\zrules} = rs}
        {$rs$ can be obtained by erasing pointer from $\zrules$}
\begin{align*}
  \rmpointer{\zruls{\cdot}{r}{rs}} &= \hrules{r}{rs} \\
  \rmpointer{\zruls{\hrulesP{r'}{rs'}}{r}{rs}} &= \hrules{r'}{\rmpointer{\zruls{rs'}{r}{rs}}}
\end{align*}

\caption{Pointer Eraser flattens the rules with pointer.}
\label{fig:pointer-eraser}
\end{figure}

% !TEX root= pattern-paper.tex

\begin{figure}[ht]

\judgbox{\htrans{e}{e'}}{$e$ takes a step to $e'$}

\begin{mathpar}
\Infer{\ITHole}{
  \htrans{e}{e'}
}{
  \htrans{\hhole{e}{u}}{\hhole{e'}{u}}
}

\Infer{\ITApFun}{
  \htrans{e_1}{e_1'}
}{
  \htrans{\hap{e_1}{e_2}}{\hap{e_1'}{e_2}}
}

\Infer{\ITApArg}{
  \isFinal{e_1} \\
  \htrans{e_2}{e_2'}
}{
  \htrans{\hap{e_1}{e_2}}{\hap{e_1}{e_2'}}
}

\Infer{\ITAp}{
  \isFinal{e_2}
}{
  \hap{\hlam{x}{\tau}{e_1}}{e_2} \mapsto
    [e_2/x]e_1
}

\Infer{\ITPairL}{
  \htrans{e_1}{e_1'}
}{
  \htrans{\hpair{e_1}{e_2}}{\hpair{e_1'}{e_2}}
}

\Infer{\ITPairR}{
  \isFinal{e_1} \\
  \htrans{e_2}{e_2'}
}{
  \htrans{\hpair{e_1}{e_2}}{\hpair{e_1}{e_2'}}
}

\Infer{\ITPrl}{
  \isFinal{\hpair{e_1}{e_2}}
}{
  \htrans{\hprl{\hpair{e_1}{e_2}}}{e_1}
}

\Infer{\ITPrr}{
  \isFinal{\hpair{e_1}{e_2}}
}{
  \htrans{\hprr{\hpair{e_1}{e_2}}}{e_2}
}

\Infer{\ITInl}{
  \htrans{e}{e'}
}{
  \htrans{\hinl{\tau}{e}}{\hinl{\tau}{e'}}
}

\Infer{\ITInr}{
  \htrans{e}{e'}
}{
  \htrans{\hinr{\tau}{e}}{\hinr{\tau}{e'}}
}

\Infer{\ITExpMatch}{
  \htrans{e}{e'}
}{
  \htrans{\hmatch{e}{\zrules}}{\hmatch{e'}{\zrules}}
}

\Infer{\ITSuccMatch}{
  \isFinal{e} \\
  \hpatmatch{e}{p_r}{\theta}
}{
  \htrans{
    \hmatch{e}{\zruls{rs_{pre}}{\hrulP{p_r}{e_r}}{rs_{post}}}
  }{
    [\theta](e_r)
  }
}

\Infer{\ITFailMatch}{
  \isFinal{e} \\
  \hnotmatch{e}{p_r}
}{
  \htrans{
    \hmatch{e}{\zruls{rs}{\hrulP{p_r}{e_r}}{\hrulesP{r'}{rs'}}}
  }{
    \hmatch{e}{
      \zruls{
        \rmpointer{\zruls{rs}{\hrulP{p_r}{e_r}}{\cdot}}
      }{r'}{rs'}
    }
  }
}
\end{mathpar}

  \caption{Stepping}
  \label{fig:step}
\end{figure}


In Sec.~\ref{sec:examples}, we use a rule pointer to simulate the process in
which we consider the rules in order. Correspondingly, in the syntax of the core
calculus, we use zipper rules of the form $\zrulsP{rs}{r}{rs}$ in match
expressions to represent the intermediate states when we are performing pattern
matching. The former $rs$ represents the preceding rules that has already been
considered and fails to match the scrutinee; the rule $r$ in the middle
represents the current rule being considered; the latter $rs$ represents the
remaining rules to be considered. \figurename~\ref{fig:pointer-eraser} defines a
helper function to flatten the zipper rules. If we need to move the rule pointer
to the next rule, we can append the current rule after the preceding rules, and
regard the initial rule in the remaining rule as the new current rule. The
conclusion of Rule \ITFailMatch demonstrates how it works.

% !TEX root= pattern-paper.tex

\begin{figure}[th]

\judgbox{\isVal{e}}{$e$ is a value}

\begin{mathpar}
\Infer{\VNum}{ }{
  \isVal{\hnum{n}}
}

\Infer{\VLam}{ }{
  \isVal{\hlam{x}{\tau}{e}}
}

\Infer{\VPair}{
  \isVal{e_1} \\
  \isVal{e_2}
}{
  \isVal{\hpair{e_1}{e_2}}
}

\Infer{\VInl}{
  \isVal{e}
}{
  \isVal{\hinl{\tau}{e}}
}

\Infer{\VInr}{
  \isVal{e}
}{
  \isVal{\hinr{\tau}{e}}
}
\end{mathpar}

\judgbox{\isIndet{e}}{$e$ is indeterminate}

\begin{mathpar}
\Infer{\IEHole}{ }{
  \isIndet{\hehole{u}}
}

\Infer{\IHole}{
  \isFinal{e}
}{
  \isIndet{\hhole{e}{u}}
}

\Infer{\IAp}{
  \isIndet{e_1} \\ \isFinal{e_2}
}{
  \isIndet{\hap{e_1}{e_2}}
}

\Infer{\IPairL}{
  \isIndet{e_1} \\ \isVal{e_2}
}{
  \isIndet{\hpair{e_1}{e_2}}
}

\Infer{\IPairR}{
  \isVal{e_1} \\ \isIndet{e_2}
}{
  \isIndet{\hpair{e_1}{e_2}}
}

\Infer{\IPair}{
  \isIndet{e_1} \\ \isIndet{e_2}
}{
  \isIndet{\hpair{e_1}{e_2}}
}

\Infer{\IFst}{
  \isIndet{e} \\ e \neq \hpair{e_1}{e_2}
}{
  \isIndet{\hfst{e}}
}

\Infer{\ISnd}{
  \isIndet{e} \\ e \neq \hpair{e_1}{e_2}
}{
  \isIndet{\hsnd{e}}
}

\Infer{\IInl}{
  \isIndet{e}
}{
  \isIndet{\hinl{\tau}{e}}
}

\Infer{\IInr}{
  \isIndet{e}
}{
  \isIndet{\hinr{\tau}{e}}
}

\Infer{\IMatch}{
  \isFinal{e} \\
  \hmaymatch{e}{p_r}
}{
  \isIndet{
    \hmatch{e}{\zruls{rs_{pre}}{\hrulP{p_r}{e_r}}{rs_{post}}}
  }
}
\end{mathpar}

\judgbox{\isFinal{e}}{$e$ is final}

\begin{mathpar}
\Infer{\FVal}{
  \isVal{e}
}{
  \isFinal{e}
}

\Infer{\FIndet}{
  \isIndet{e}
}{
  \isFinal{e}
}
\end{mathpar}

  \caption{Final expressions}
  \label{fig:final}
\end{figure}


\figurename~\ref{fig:step} defines the stepping judgment
$\htrans{e}{e'}$. We will focus on stepping judgments of match expressions.
Rule \ITExpMatch specifies that the match expression can take a step as its its
scrutinee can take a step. Note that Rule \ITFailMatch and Rule \ITSuccMatch
share the premise, $\isFinal{e}$, which is defined in \figurename~\ref{fig:final}. It
means that expression $e$ is \textit{final} in the sense that it is either already a
\textit{value} or \textit{indeterminate}, \ie, cannot be evaluated further due to unfilled holes
\cite{DBLP:journals/pacmpl/OmarVCH19}. And only when the scrutinee is final,
shall we consider the constinuent rules of the match expression in order. As in
the examples shown in \listfigurename~\ref{fig:exp-hole-step},\ref{fig:pat-hole-step},

\begin{itemize}
  \item
    if the final scrutinee does not match the pattern in the current rule,
    then we move the pointer to the next rule (Rule \ITFailMatch)

  \item
    if the final scrutinee does match the pattern in the current rule, 
    then we take the emitted substitution $\theta$ and apply it on the subexpression of that rule (Rule \ITSuccMatch)

  \item 
    if the final scrutinee may or may not match the pattern in the current rule,
    then the match expression is said to be indeterminate (Rule \IMatch)
\end{itemize}

% !TEX root= pattern-paper.tex

\begin{figure}[!b]

\judgbox{
  \hpatmatch{e}{p}{\theta}
}{
  $e$ matches $p$, emitting $\theta$
}

\begin{mathpar}
\Infer{\MVar}{ }{
  \hpatmatch{e}{x}{e / x}
}

\Infer{\MWild}{ }{
  \hpatmatch{e}{\_}{\cdot}
}

\Infer{\MNum}{ }{
  \hpatmatch{\hnum{n}}{\hnum{n}}{\cdot}
}

\Infer{\MPair}{
  \hpatmatch{e_1}{p_1}{\theta_1} \\
  \hpatmatch{e_2}{p_2}{\theta_2}
}{
  \hpatmatch{\hpair{e_1}{e_2}}{\hpair{p_1}{p_2}}{\theta_1 \uplus \theta_2}
}

\Infer{\MInl}{
  \hpatmatch{e}{p}{\theta}
}{
  \hpatmatch{\hinl{\tau}{e}}{\hinlp{p}}{\theta}
}

\Infer{\MInr}{
  \hpatmatch{e}{p}{\theta}
}{
  \hpatmatch{\hinr{\tau}{e}}{\hinrp{p}}{\theta}
}

\Infer{\MNotIntroPair}{
  \notIntro{e} \\
  \hpatmatch{\hprl{e}}{p_1}{\theta_1} \\
  \hpatmatch{\hprr{e}}{p_2}{\theta_2}
}{
  \hpatmatch{e}{\hpair{p_1}{p_2}}{\theta_1 \uplus \theta_2}
}
\end{mathpar}

\judgbox{
  \hnotmatch{e}{p}
}{
  $e$ does not match $p$
}

\begin{mathpar}

\Infer{\NMNum}{
  n_1 \neq n_2
}{
  \hnotmatch{\hnum{n_1}}{\hnum{n_2}}
}

\Infer{\NMPairL}{
  \hnotmatch{e_1}{p_1}
}{
  \hnotmatch{\hpair{e_1}{e_2}}{\hpair{p_1}{p_2}}
}

\Infer{\NMPairR}{
  \hnotmatch{e_2}{p_2}
}{
  \hnotmatch{\hpair{e_1}{e_2}}{\hpair{p_1}{p_2}}
}

\Infer{\NMConfL}{ }{
  \hnotmatch{\hinr{\tau}{e}}{\hinlp{p}}
}

\Infer{\NMConfR}{ }{
  \hnotmatch{\hinl{\tau}{e}}{\hinrp{p}}
}

\Infer{\NMInl}{
  \hnotmatch{e}{p}
}{
  \hnotmatch{\hinr{\tau}{e}}{\hinlp{p}}
}

\Infer{\NMInr}{
  \hnotmatch{e}{p}
}{
  \hnotmatch{\hinl{\tau}{e}}{\hinrp{p}}
}
\end{mathpar}

\judgbox{
  \hmaymatch{e}{p}
}{
  $e$ indeterminately matches $p$
}

\begin{mathpar}
\Infer{\MMEHole}{ }{
  \hmaymatch{e}{\heholep{w}}
}

\Infer{\MMHole}{ }{
  \hmaymatch{e}{\hholep{p}{w}{\tau}}
}
\\
\Infer{\MMNotIntro}{
  \notIntro{e} \\
  \refutable{p}
}{
  \hmaymatch{e}{p}
}

\Infer{\MMPairL}{
  \hmaymatch{e_1}{p_1} \\
  \hpatmatch{e_2}{p_2}{\theta_2}
}{
  \hmaymatch{\hpair{e_1}{e_2}}{\hpair{p_1}{p_2}}
}

\Infer{\MMPairR}{
  \hpatmatch{e_1}{p_1}{\theta_1} \\
  \hmaymatch{e_2}{p_2}
}{
  \hmaymatch{\hpair{e_1}{e_2}}{\hpair{p_1}{p_2}}
}

\Infer{\MMPair}{
  \hmaymatch{e_1}{p_1} \\
  \hmaymatch{e_2}{p_2}
}{
  \hmaymatch{\hpair{e_1}{e_2}}{\hpair{p_1}{p_2}}
}

\Infer{\MMInl}{
  \hmaymatch{e}{p}
}{
  \hmaymatch{\hinl{\tau}{e}}{\hinlp{p}}
}

\Infer{\MMInr}{
  \hmaymatch{e}{p}
}{
  \hmaymatch{\hinr{\tau}{e}}{\hinrp{p}}
}
\end{mathpar}

\caption{Three possible outcomes of pattern matching}
\label{fig:patmatch}
\end{figure}


\figurename~\ref{fig:patmatch} defines three judgments corresponding to the
three possible results of pattern matching. We consider a final scrutinee $e$
and a pattern $p$ that are of the same type (see Sec.~\ref{sec:statics}). The
judgment $\hpatmatch{e}{p}{\theta}$ denotes that $e$ matches $p$, as witness by
the substitution $\theta$ defined on the variables in $p$, while the judgment
$\hnotmatch{e}{p}$ denotes that $e$ doesn't match $p$. Notably, we introduce the
third possibility $\hmaymatch{e}{p}$ --- $e$ may match $p$, which is also the
\textit{key concept} in pattern matching with holes. Rules \MMEHole and \MMHole
specifies that any final expression may match a pattern hole, regardless of
whether it is empty or non-empty. Rules \MMExpEHole and \MMExpHole specifies
that an expression hole, regardless of whether it is empty or non-empty, may
match any pattern except variable patterns and wild card patterns. It is because
Rules \MVar and \MWild specifies any final expression satisfies a variable
pattern or a wildcard pattern, and it does not make sense to allow
$\hpatmatch{e}{p}{\theta}$ and $\hmaymatch{e}{p}$ at the same time. Actually,
the three possible results of pattern matching are mutually exclusive and at
least one of judgments can be derived under certain assumptions, as shown in the
following lemma.

\begin{lemma}[Matching Determinism]
  \label{lem:match-determinism}
  If $\isFinal{e}$ and $\hexptyp{\cdot}{\Delta}{e}{\tau}$ and $\hpattyp{p}{\tau}{\Gamma}{\Delta}$ then exactly one of the following holds
  \begin{enumerate}
    \item $\hpatmatch{e}{p}{\theta}$ for some $\theta$
    \item $\hmaymatch{e}{p}$
    \item $\hnotmatch{e}{p}$
  \end{enumerate}
\end{lemma}

The determinism of pattern matching ensures the determinism of dynamic semantics.

\begin{theorem}[Determinism]
  \label{thrm:determinism}
  If $\hexptyp{\cdot}{\Delta}{e}{\tau}$ then exactly one of the following holds
  \begin{enumerate}
    \item $\isVal{e}$
    \item $\isIndet{e}$
    \item $\htrans{e}{e'}$ for some unique $e'$
  \end{enumerate}
\end{theorem}

\todo{State progress theorem instead of determinism or both?}

\begin{lemma}[Matching Progress]
  \label{lem:match-progress}
  If $\isFinal{e}$ and $\hexptyp{\cdot}{\Delta}{e}{\tau}$ and $\hpattyp{p}{\tau}{\Gamma}{\Delta'}$
  then $\hnotmatch{e}{p}$ or $\hmaymatch{e}{p}$ or $\hpatmatch{e}{p}{\theta}$ for some $\theta$
\end{lemma}

\begin{theorem}[Progress]
  \label{thrm:progress}
  If $\hexptyp{\cdot}{\Delta}{e}{\tau}$ then either $\isFinal{e}$ or $\htrans{e}{e'}$ for some $e'$.
\end{theorem}

\subsection{Static Semantics}\label{sec:statics}

We have already introduced how pattern matching with holes works. However, in order to do exhaustiveness checking and redundancy
checking of match expressions in \listfigurename~\ref{fig:exh-hole},
\ref{fig:red-hole}, we need to predict the result of pattern matching in static
semantics. We start by introducing \textit{match constraint language}, which
extends the idea in \cite{Harper2012}. Then, we build a similar type system to
\cite{DBLP:journals/pacmpl/OmarVCH19} by defining typing judgments for both
patterns and expressions. The former generates variable contexts $\Gamma$ and
hole contexts $\Delta$ (see \figurename~\ref{fig:pattyp}) while the latter 
takes variable contexts $\Gamma$ and hole contexts $\Delta$ as hypothesis (see
\figurename~\ref{fig:exptyp}).

\subsubsection{Match Constraint Language}\label{sec:constraint}
% !TEX root= pattern-paper.tex

\begin{figure}[ht]
$\arraycolsep=4pt\begin{array}{lll}
\xi & ::= &
  \ctruth ~\vert~
  \cfalsity ~\vert~
  \cunknown ~\vert~
  \cnum{n} ~\vert~
  \cnotnum{n} ~\vert~
  \cand{\xi_1}{\xi_2} ~\vert~
  \cor{\xi_1}{\xi_2} ~\vert~
  \cinl{\xi} ~\vert~
  \cinr{\xi} ~\vert~
  \cpair{\xi_1}{\xi_2}
\end{array}$

\judgbox{\ctyp{\xi}{\tau}}{$\xi$ constrains values of type $\tau$}

\begin{mathpar}
\Infer{\CTTruth}{ }{
  \ctyp{\ctruth}{\tau}
}

\Infer{\CTFalsity}{ }{
  \ctyp{\cfalsity}{\tau}
}

\Infer{\CTUnknown}{ }{
  \ctyp{\cunknown}{\tau}
}

\Infer{\CTNum}{ }{
  \ctyp{\cnum{n}}{\tnum}
}

\Infer{\CTNotNum}{ }{
  \ctyp{\cnotnum{n}}{\tnum}
}

\Infer{\CTAnd}{
  \ctyp{\xi_1}{\tau} \\ \ctyp{\xi_2}{\tau}
}{
  \ctyp{\cand{\xi_1}{\xi_2}}{\tau}
}

\Infer{\CTOr}{
  \ctyp{\xi_1}{\tau} \\ \ctyp{\xi_2}{\tau}
}{
  \ctyp{\cor{\xi_1}{\xi_2}}{\tau}
}

\Infer{\CTInl}{
  \ctyp{\xi_1}{\tau_1}
}{
  \ctyp{\cinl{\xi_1}}{\tsum{\tau_1}{\tau_2}}
}

\Infer{\CTInr}{
  \ctyp{\xi_2}{\tau_2}
}{
  \ctyp{\cinr{\xi_2}}{\tsum{\tau_1}{\tau_2}}
}

\Infer{\CTPair}{
  \ctyp{\xi_1}{\tau} \\ \ctyp{\xi_2}{\tau}
}{
  \ctyp{\cpair{\xi_1}{\xi_2}}{\tau}
}
\end{mathpar}

\judgbox{\cdual{\xi_1} = \xi_2}{dual of $\xi_1$ is $\xi_2$}
\begin{subequations}\label{defn:dual}
\begin{align}
  \cdual{\ctruth} &= \cfalsity \\
  \cdual{\cfalsity} &= \ctruth \\
  \cdual{\cunknown} &= \cunknown \\
  \cdual{\cnum{n}} &= \cnotnum{n} \\
  \cdual{\cnotnum{n}} &= \cnum{n} \\
  \cdual{\cand{\xi_1}{\xi_2}} &= \cor{\cdual{\xi_1}}{\cdual{\xi_2}} \\
  \cdual{\cor{\xi_1}{\xi_2}} &= \cand{\cdual{\xi_1}}{\cdual{\xi_2}} \\
  \cdual{\cinl{\xi_1}} &= \cor{ \cinl{\cdual{\xi_1}} }{ \cinr{\ctruth} } \\
  \cdual{\cinr{\xi_2}} &= \cor{ \cinr{\cdual{\xi_2}} }{ \cinl{\ctruth} } \\
  \cdual{\cpair{\xi_1}{\xi_2}} &=
  \cor{ \cor{ 
    \cpair{\xi_1}{\cdual{\xi_2}}
  }{
    \cpair{\cdual{\xi_1}}{\xi_2}
  }}{
    \cpair{\cdual{\xi_1}}{\cdual{\xi_2}}
  }
\end{align}
\end{subequations}
  \caption{Match Constraints}
  \label{fig:constraint}
\end{figure}


\figurename~\ref{fig:constraint} gives the syntax of match constraint, which is
used to identify a subset of the final expressions of a type. The judgment
$\ctyp{\xi}{\tau}$ specifies that constraint $\xi$ constrains the final
expressions of type $\tau$. The dual of $\xi$, $\cdual{\xi}$ represents the
complement of the subset identified by $\xi$ under the set of the final
expressions of a type.

\subsubsection{Typing of Patterns}
% !TEX root= extended-abstract.tex

\begin{figure}[h]
\judgbox{\chpattyp{p}{\tau}{\xi}{\Gamma}{\Delta}}
{$p$ matches expression of type $\tau$, with match constraint $\xi$}

\begin{mathpar}
\Infer{\PTVar}{ }{
  \chpattyp{x}{\tau}{\ctruth}{\cdot}{x : \tau}
}

\Infer{\PTNum}{ }{
  \chpattyp{\hnum{n}}{\tnum}{\cnum{n}}{\cdot}{\cdot}
}

\Infer{\PTWild}{ }{
  \chpattyp{\_}{\tau}{\ctruth}{\cdot}{\cdot}
}

\Infer{\PTInL}{
  \chpattyp{p}{\tau_1}{\xi}{\Gamma}{\Delta}
}{
  \chpattyp{\hinlp{p}}{\tsum{\tau_1}{\tau_2}}{\cinl{\xi}}{\Gamma}{\Delta}
}

\Infer{\PTInR}{
  \chpattyp{p}{\tau_2}{\xi}{\Gamma}{\Delta}
}{
  \chpattyp{\hinrp{p}}{\tsum{\tau_1}{\tau_2}}{\cinr{\xi}}{\Gamma}{\Delta}
}

\Infer{\PTPair}{
  \chpattyp{p_1}{\tau_1}{\xi_1}{\Gamma_1}{\Delta_1} \\
  \chpattyp{p_2}{\tau_2}{\xi_2}{\Gamma_2}{\Delta_2}
}{
  \chpattyp{\hpair{p_1}{p_2}}{\tprod{\tau_1}{\tau_2}}
  {\cpair{\xi_1}{\xi_2}}{\Gamma_1 \uplus \Gamma_2}{\Delta_1 \uplus \Delta_2}
}

\Infer{\PTEHole}{ }{
  \chpattyp{\hehole{w}}{\tau}{\cunknown}{\cdot}{w :: \tau}
}

\Infer{\PTHole}{
  \chpattyp{p}{\tau}{\xi}{\Gamma}{\Delta}
}{
  \chpattyp{\hhole{p}{w}}{\tau'}{\cunknown}
  {\Gamma}{\Delta , w :: \tau'}
}
\end{mathpar}
\caption{Typing of patterns}
\label{fig:pattyp}
\end{figure}


\figurename~\ref{fig:pattyp} defines the typing judgment for patterns and
connects constraints with patterns. When pattern $p$ is assigned type $\tau$,
we collect the typing for the pattern variables and pattern holes in $p$ and
generate the corresponding variable context $\Gamma$ and hole context $\Delta$.
Besides two contexts, we emit a constraint $\xi$ to identify the set of final
expressions that match $p$. Rules \PTEHole and Rules \PTHole specifies that
pattern holes emits \textit{Unknown} constraint, $\cunknown$. That means, the
set of final expressions that matches the pattern is yet to be identified, we
will see how it is used in \figurename~\ref{fig:satisfy}.

\subsubsection{Satisfaction Judgments}
% !TEX root= pattern-paper.tex

\begin{figure}[!ht]

\judgbox{\csatisfy{e}{\xi}}{$e$ satisfies $\xi$}

\begin{mathpar}
\Infer{\CSTruth}{ }{
  \csatisfy{e}{\ctruth}
}

\Infer{\CSNum}{ }{
  \csatisfy{\hnum{n}}{\cnum{n}}
}

\Infer{\CSNotNum}{
  n_1 \neq n_2
}{
  \csatisfy{\hnum{n_1}}{\cnotnum{n_2}}
}

\Infer{\CSAnd}{
  \csatisfy{e}{\xi_1} \\
  \csatisfy{e}{\xi_2}
}{
  \csatisfy{e}{\cand{\xi_1}{\xi_2}}
}

\Infer{\CSOrL}{
  \csatisfy{e}{\xi_1}
}{
  \csatisfy{e}{\cor{\xi_1}{\xi_2}}
}

\Infer{\CSOrR}{
  \csatisfy{e}{\xi_2}
}{
  \csatisfy{e}{\cor{\xi_1}{\xi_2}}
}

\Infer{\CSInl}{
  \csatisfy{e_1}{\xi_1}
}{
  \csatisfy{
    \hinl{\tau_2}{e_1}
  }{
    \cinl{\xi_1}
  }
}

\Infer{\CSInr}{
  \csatisfy{e_2}{\xi_2}
}{
  \csatisfy{
    \hinr{\tau_1}{e_2}
  }{
    \cinr{\xi_2}
  }
}

\Infer{\CSPair}{
  \csatisfy{e_1}{\xi_1} \\
  \csatisfy{e_2}{\xi_2}
}{
\csatisfy{\hpair{e_1}{e_2}}{\cpair{\xi_1}{\xi_2}}
}

\Infer{\CSNotIntroPair}{
  \notIntro{e} \\
  \csatisfy{\hprl{e}}{\xi_1} \\
  \csatisfy{\hprr{e}}{\xi_2}
}{
  \csatisfy{e}{\cpair{\xi_1}{\xi_2}}
}
\end{mathpar}

\judgbox{\cmaysatisfy{e}{\xi}}{$e$ may satisfy $\xi$}

\begin{mathpar}
\Infer{\CMSUnknown}{ }{
  \cmaysatisfy{e}{\cunknown}
}

\Infer{\CMSNotIntro}{
  \notIntro{e} \\
  \refutable{\xi}
}{
  \cmaysatisfy{e}{\xi}
}

\Infer{\CMSAndL}{
  \cmaysatisfy{e}{\xi_1} \\
  \csatisfy{e}{\xi_2}
}{
  \cmaysatisfy{e}{\cand{\xi_1}{\xi_2}}
}

\Infer{\CMSAndR}{
  \csatisfy{e}{\xi_1} \\
  \cmaysatisfy{e}{\xi_2}
}{
  \cmaysatisfy{e}{\cand{\xi_1}{\xi_2}}
}

\Infer{\CMSAnd}{
  \cmaysatisfy{e}{\xi_1} \\
  \cmaysatisfy{e}{\xi_2}
}{
  \cmaysatisfy{e}{\cand{\xi_1}{\xi_2}}
}

\Infer{\CMSOrL}{
  \cmaysatisfy{e}{\xi_1} \\
  \cnotsatisfy{e}{\xi_2}
}{
  \cmaysatisfy{e}{\cor{\xi_1}{\xi_2}}
}

\Infer{\CMSOrR}{
  \cnotsatisfy{e}{\xi_1} \\
  \cmaysatisfy{e}{\xi_2}
}{
  \cmaysatisfy{e}{\cor{\xi_1}{\xi_2}}
}

\Infer{\CMSInl}{
  \cmaysatisfy{e_1}{\xi_1}
}{
  \cmaysatisfy{
    \hinl{\tau_2}{e_1}
  }{
    \cinl{\xi_1}
  }
}

\Infer{\CMSInr}{
  \cmaysatisfy{e_2}{\xi_2}
}{
  \cmaysatisfy{
    \hinr{\tau_1}{e_2}
  }{
    \cinr{\xi_2}
  }
}

\Infer{\CMSPairL}{
  \cmaysatisfy{e_1}{\xi_1} \\
  \csatisfy{e_2}{\xi_2}
}{
  \cmaysatisfy{\hpair{e_1}{e_2}}{\cpair{\xi_1}{\xi_2}}
}

\Infer{\CMSPairR}{
  \csatisfy{e_1}{\xi_1} \\
  \cmaysatisfy{e_2}{\xi_2}
}{
  \cmaysatisfy{\hpair{e_1}{e_2}}{\cpair{\xi_1}{\xi_2}}
}

\Infer{\CMSPair}{
  \cmaysatisfy{e_1}{\xi_1} \\
  \cmaysatisfy{e_2}{\xi_2}
}{
  \cmaysatisfy{\hpair{e_1}{e_2}}{\cpair{\xi_1}{\xi_2}}
}
\end{mathpar}

\judgbox{\csatisfyormay{e}{\xi}}{$e$ satisfies or may satisfy $\xi$}

\begin{mathpar}
\Infer{\CSMSMay}{
  \cmaysatisfy{e}{\xi}
}{
  \csatisfyormay{e}{\xi}
}

\Infer{\CSMSSat}{
  \csatisfy{e}{\xi}
}{
  \csatisfyormay{e}{\xi}
}
\end{mathpar}

  \caption{Satisfaction Judgments}
  \label{fig:satisfy}
\end{figure}


\figurename~\ref{fig:satisfy} defines satisfaction judgments. As we only
consider final expressions and patterns of the same type when talking about
pattern matching, a constraint only constrains final expressions of the same
type. And the satisfaction judgments does not make sense when the expression is
not final or the expression and the constraint is of different type. The
judgment $\csatisfy{e}{\xi}$ specifies that expression $e$ satisfies $\xi$ while
the judgment $\cmaysatisfy{e}{\xi}$ specifies that expression $e$ may or may not
satisfy $\xi$. The judgment $\csatisfyormay{e}{\xi}$ is the combination of the
previous two cases. It turns out that the remaining case where
$\csatisfyormay{e}{\xi}$ is not derivable, can be represented by
$\csatisfy{e}{\cdual{\xi}}$.

\begin{theorem}[Exclusiveness of Constraint Satisfaction]
  \label{thrm:exclusive-constraint-satisfaction}
  If $\ctyp{\xi}{\tau}$ and $\hexptyp{\cdot}{\Delta}{e}{\tau}$ and $\isFinal{e}$ then exactly one of the following holds
  \begin{enumerate}
  \item $\csatisfy{e}{\xi}$
  \item $\cmaysatisfy{e}{\xi}$
  \item $\csatisfy{e}{\cdual{\xi}}$
  \end{enumerate}
\end{theorem}

\lemmaname~\ref{lem:const-matching-coherence} establishes a correspondence
between pattern matching results and satisfaction judgments. That makes
reasoning pattern matching in type system possible and helps prove
\theoremname~\ref{thrm:preservation}.

\begin{lemma}[Matching Coherence of Constraint]
  \label{lem:const-matching-coherence}
  Suppose that $\hexptyp{\cdot}{\Delta}{e}{\tau}$ and $\isFinal{e}$ and $\chpattyp{p}{\tau}{\xi}{\Gamma}{\Delta}$. Then we have
  \begin{enumerate}
  \item $\csatisfy{e}{\xi}$ iff $\hpatmatch{e}{p}{\theta}$
  \item $\csatisfy{e}{\cdual{\xi}}$ iff $\hnotmatch{e}{p}$
  \item $\cmaysatisfy{e}{\xi}$ iff $\hmaymatch{e}{p}$
  \end{enumerate}
\end{lemma}

The following two definitions take advantage satisfaction judgments and
corresponds to redundancy and exhaustiveness respectively. We will see how they
can be determined in Sec.~\ref{sec:algorithm}

\begin{definition}[Entailment of Constraints]
  \label{defn:const-entailment}
  Suppose that $\ctyp{\xi_1}{\tau}$ and $\ctyp{\xi_2}{\tau}$.
  Then $\csatisfy{\xi_1}{\xi_2}$ iff for all $e$ such that $\hexptyp{\cdot}{\Delta}{e}{\tau}$ and $\isFinal{e}$ we have $\csatisfyormay{e}{\xi_1}$ implies $\csatisfy{e}{\xi_2}$
\end{definition}

\begin{definition}[Not Not Entailment of Constraints]
  \label{defn:nn-entailment}
  Suppose that $\ctyp{\xi_1}{\tau}$ and $\ctyp{\xi_2}{\tau}$. Then $\csatisfyormay{\xi_1}{\xi_2}$ iff for all $e$ such that $\hexptyp{\cdot}{\Delta}{e}{\tau}$ and $\isFinal{e}$ we have $\csatisfyormay{e}{\xi_1}$ implies $\csatisfyormay{e}{\xi_2}$ 
\end{definition}

\subsubsection{Typing of Rules and Redundancy Checking}

% !TEX root= extended-abstract.tex

\begin{figure}[h]
\judgbox{\chrultyp{\Gamma}{\Delta}{r}{\tau_1}{\xi}{\tau_2}}
{$r$ is of type $\trulnp{\tau_1}{\tau_2}$ and is constrained by $\xi$}

\begin{mathpar}
\Infer{\TRule}{
  \chpattyp{p}{\tau_1}{\xi}{\Gamma_p}{\Delta_p} \\
  \hexptyp{\Gamma \uplus \Gamma_p}{\Delta \uplus \Delta_p}{e}{\tau_2}
}{
  \chrultyp{\Gamma}{\Delta}{\hrul{p}{e}}{\tau_1}{\xi}{\tau_2}
}
\end{mathpar}
\caption{Typing of rule}
\label{fig:ruletyp}
\end{figure}
 

\figurename~\ref{fig:ruletyp} defines the typing judgment of a rule. For a rule
$\hrul{p}{e}$, after generating the variable context and hole context from the
typing of $p$, we extend the variable context and the hole context of the rule
with the generated contexts and form the hypothesis for the typing of its
subexpression. Also, a rule inherits the constraint from its pattern.

% !TEX root= extended-abstract.tex

\begin{figure}[h]

\begin{mathpar}
\Infer{\TRules}{
  \chrultyp{\Gamma}{\Delta}{r}{\tau_1}{\xi_r}{\tau_2} \\
  \chrulstyp{\Gamma}{\Delta}{\cor{\xi_{pre}}{\xi_r}}{rs}
  {\tau_1}{\xi_{rs}}{\tau_2} \\
  \boxed{\cnotsatisfy{\xi_r}{\xi_{pre}}}
}{
  \chrulstyp{\Gamma}{\Delta}{\xi_{pre}}{\hrules{r}{rs}}
  {\tau_1}{\cor{\xi_r}{\xi_{rs}}}{\tau_2}
}
\end{mathpar}
\label{fig:rulestyp}
\end{figure}


\figurename~\ref{fig:rulestyp} defines the typing judgment of a series of rules
that at least has one rule. Compared to a single rule, the form of the typing
judgment of rules is slightly different and involves two constraint components,
where $\xi_{pre}$ represent the constraints emitted from its preceding rules and
$\xi_{rs}$ represents the constraint emitted from the series of rules being type
checked. As we build up the rules, we check if the initial rule is redundant
with respect to its preceding rules, and the premise in both two Rules,
$\cnotsatisfy{\xi_r}{\xi_{pre}}$. It states that the judgment
$\csatisfy{\xi_r}{\xi_{pre}}$, which corresponds to redundancy, is not
derivable. In other words, rule $r$ either cannot be redundant or, may or may
not be redundant, with respect to it preceding rules. Rule \TOneRules specifies
that for \todo{single rules}, we can simply inherit $\xi_r$ emitted from typing
of $r$ and check the redundancy of $r$. Rule \TRules specifies that we collect
the $\xi_r$ and concatenate it with $\xi_{pre}$ as the new input to the typing
judgment of $rs$ and then wait for the output $\xi_rs$ and concatenate it with
$\xi_r$ to complete the conclusion.

\subsubsection{Typing of Expressions and Exhaustiveness Checking}

We have already shown how to collect constraints from rules in the typing
judgment of rules. Now we can use collected constraints to check exhaustiveness
in the typing judgment of match expressions.

% !TEX root= pattern-paper.tex

\begin{figure}[ht]
\judgbox{
  \hexptyp{\Gamma}{\Delta}{e}{\tau}
}{
  $e$ is of type \(\tau\)
}
  \begin{mathpar}
%  \Infer{\TVar}{ }{
%    \hexptyp{\Gamma, x:\tau}{\Delta}{x}{\tau}
%  }
%
%  \Infer{\TEHole}{ }{
%    \hexptyp{\Gamma}{\Delta, u::\tau}{\hehole{u}}{\tau}
%  }
%
%  \Infer{\THole}{
%    \hexptyp{\Gamma}{\Delta, u::\tau}{e}{\tau'}
%  }{
%    \hexptyp{\Gamma}{\Delta, u::\tau}{\hhole{e}{u}}{\tau}
%  }
%
%  \Infer{\TNum}{ }{
%    \hexptyp{\Gamma}{\Delta}{\hnum{n}}{\tnum}
%  }
%
%  \Infer{\TLam}{
%    \hexptyp{\Gamma, x:\tau_1}{\Delta}{e}{\tau_2}
%  }{
%    \hexptyp{\Gamma}{\Delta}{\hlam{x}{\tau_1}{e}}{\tarr{\tau_1}{\tau_2}}
%  }
%
%  \Infer{\TAp}{
%    \hexptyp{\Gamma}{\Delta}{e_1}{\tarr{\tau_2}{\tau}} \\
%    \hexptyp{\Gamma}{\Delta}{e_2}{\tau_2}
%  }{
%    \hexptyp{\Gamma}{\Delta}{\hap{e_1}{e_2}}{\tau}
%  }
%
%  \Infer{\TPair}{
%    \hexptyp{\Gamma}{\Delta}{e_1}{\tau_1} \\
%    \hexptyp{\Gamma}{\Delta}{e_2}{\tau_2}
%  }{
%    \hexptyp{\Gamma}{\Delta}{\hpair{e_1}{e_2}}{\tprod{\tau_1}{\tau_2}}
%  }
%
%  \Infer{\TFst}{
%    \hexptyp{\Gamma}{\Delta}{e}{\tprod{\tau_1}{\tau_2}}
%  }{
%    \hexptyp{\Gamma}{\Delta}{\hfst{e}}{\tau_1}
%  }
%  
%  \Infer{\TSnd}{
%    \hexptyp{\Gamma}{\Delta}{e}{\tprod{\tau_1}{\tau_2}}
%  }{
%    \hexptyp{\Gamma}{\Delta}{\hsnd{e}}{\tau_2}
%  }
%  
%  \Infer{\TInl}{
%    \hexptyp{\Gamma}{\Delta}{e}{\tau_1}
%  }{
%    \hexptyp{\Gamma}{\Delta}{\hinl{\tau_2}{e}}{\tsum{\tau_1}{\tau_2}}
%  }
%
%  \Infer{\TInr}{
%    \hexptyp{\Gamma}{\Delta}{e}{\tau_2}
%  }{
%    \hexptyp{\Gamma}{\Delta}{\hinr{\tau_1}{e}}{\tsum{\tau_1}{\tau_2}}
%  }
  
  \Infer{TMatchZPre}{
    \hexptyp{\Gamma}{\Delta}{e}{\tau} \\
    \chrulstyp{\Gamma}{\Delta}{\cfalsity}{\hrules{r}{rs}}{\tau}{\xi}{\tau'} \\
    \csatisfyormay{\ctruth}{\xi}
  }{
  \hexptyp{\Gamma}{\Delta}{\hmatch{e}{\zruls{\cdot}{r}{rs}}}{\tau'}
  }

  \Infer{TMatchNZPre}{
    \hexptyp{\Gamma}{\Delta}{e}{\tau} \\
    \isFinal{e} \\
    \chrulstyp{\Gamma}{\Delta}{\cfalsity}{rs_{pre}}{\tau}{\xi_{pre}}{\tau'} \\
    \chrulstyp{\Gamma}{\Delta}{\xi_{pre}}{\hrules{r}{rs_{post}}}{\tau}{\xi_{rest}}{\tau'} \\
    \cnotsatisfyormay{e}{\xi_{pre}} \\
    \csatisfyormay{\ctruth}{\cor{\xi_{pre}}{\xi_{rest}}}
  }{
    \hexptyp{\Gamma}{\Delta}{\hmatch{e}{\zruls{rs_{pre}}{r}{rs_{post}}}}{\tau'}
  }
  \end{mathpar}

\caption{Match Expression Typing}
\label{fig:exptyp}
\end{figure}


\figurename~\ref{fig:exptyp} defines the typing judgment of expressions. Rule
\TMatchZPre corresponds to the case where we haven't started pattern matching.
As shown in \definitionname~\ref{defn:nn-entailment}, the premise
$\csatisfyormay{\ctruth}{\xi}$ ensures that there is at least one rule whose
pattern does match or may match the final scrutinee and thus the pattern
matching either results in a successful match (Rule \ITSuccMatch) or an
indeterminate match expression (Rule \IMatch). Only when the final scrutinee $e$
cannot match the pattern $p$, \ie, $\hnotmatch{e}{p}$, can we move the rule
pointer. By \lemmaname~\ref{lem:match-determinism}, for all patterns $p$ in the
preceding rules, neither $\hpatmatch{e}{p}{\theta}$ nor $\hmaymatch{e}{p}$ is
derivable. Then by \lemmaname~\ref{lem:const-matching-coherence}, we can derive the
premise in Rule \TMatchNZPre, $\cnotsatisfyormay{e}{\xi_{pre}}$. And thus, the
type of the match expression is preserved as we consider rules in order.

\begin{theorem}[Preservation]
  \label{thrm:preservation}
  If $\hexptyp{\cdot}{\Delta}{e}{\tau}$ and $\htrans{e}{e'}$
  then $\hexptyp{\cdot}{\Delta}{e'}{\tau}$
\end{theorem}

% !TEX root= pattern-paper.tex

\begin{figure}[ht]

\judgbox{
  \hsubstyp{\theta}{\Gamma}
}{
  $\theta$ is of type $\Gamma$
}

\begin{mathpar}
\Infer{\STEmpty}{ }{
  \hsubstyp{\emptyset}{\cdot}
}

\Infer{\STExtend}{
  \hsubstyp{\theta}{\Gamma_\theta} \\
  \hexptyp{\Gamma}{\Delta}{e}{\tau}
}{
  \hsubstyp{\theta , x / e}{\Gamma_\theta , x : \tau}
}
\end{mathpar}

\caption{Typing of Substitution}
\label{fig:substyp}
\end{figure}


\figurename~\ref{fig:substyp} defines the typing of substitution $\theta$.

To prove \theoremname~\ref{thrm:preservation}, we need the following three lemmas.
When considering Rule \ITAp, \lemmaname~\ref{lem:substitution} is needed.
When considering Rule \ITSuccMatch, \lemmaname~\ref{lem:subs-typing} is needed
to show the typing of substitution $\theta$, and we use
\lemmaname~\ref{lem:simult-substitution} to show that type is preserved when pattern
matching succeeds.

\begin{lemma}[Substitution]
  \label{lem:substitution}
  If $\hexptyp{\Gamma, x : \tau}{\Delta}{e_0}{\tau_0}$ and $\hexptyp{\Gamma}{\Delta}{e}{\tau}$
  then $\hexptyp{\Gamma}{\Delta}{[e/x]e_0}{\tau_0}$
\end{lemma}

\begin{lemma}[Substitution Typing]
  \label{lem:subs-typing}
  If $\hpatmatch{e}{p}{\theta}$ and $\hexptyp{\cdot}{\Delta}{e}{\tau}$ and $\hpattyp{p}{\tau}{\Gamma}{\Delta'}$
  then $\hsubstyp{\theta}{\Gamma}$
\end{lemma}

\begin{lemma}[Simultaneous Substitution]
  \label{lem:simult-substitution}
  If $\hexptyp{\Gamma \uplus \Gamma'}{\Delta}{e}{\tau}$ and $\hsubstyp{\theta}{\Gamma'}$
  then $\hexptyp{\Gamma}{\Delta}{[\theta]e}{\tau}$
\end{lemma}

\section{De-unknown}\label{sec:algorithm}

We have already shown in Sec.~\ref{sec:statics} how to check redundancy and
exhaustiveness using \definitionname~\ref{defn:const-entailment} and
\definitionname~\ref{defn:nn-entailment}, but it remains unclear how to
determine whether they are true or false.

% !TEX root= pattern-paper.tex

\begin{figure}[ht]
\judgbox{\ctruify{\xi_1} = \xi_2}{}

\begin{align*}
  \ctruify{\ctruth} &= \ctruth \\
  \ctruify{\cfalsity} &= \cfalsity \\
  \ctruify{\cunknown} &= \ctruth \\
  \ctruify{\cnum{n}} &= \cnum{n} \\
  \ctruify{\cnotnum{n}} &= \cnotnum{n} \\
  \ctruify{\cand{\xi_1}{\xi_2}} &= \cand{\ctruify{\xi_1}}{\ctruify{\xi_2}} \\
  \ctruify{\cor{\xi_1}{\xi_2}} &= \cor{\ctruify{\xi_1}}{\ctruify{\xi_2}} \\
  \ctruify{\cinl{\xi}} &= \cinl{\ctruify{\xi}} \\
  \ctruify{\cinr{\xi}} &= \cinr{\ctruify{\xi}} \\
  \ctruify{\cpair{\xi_1}{\xi_2}} &= \cpair{\ctruify{\xi_1}}{\ctruify{\xi_2}}
\end{align*}

\judgbox{\cfalsify{\xi_1} = \xi_2}{}
\begin{align*}
  \cfalsify{\ctruth} &= \ctruth \\
  \cfalsify{\cfalsity} &= \cfalsity \\
  \cfalsify{\cunknown} &= \cfalsity \\
  \cfalsify{\cnum{n}} &= \cnum{n} \\
  \cfalsify{\cnotnum{n}} &= \cnotnum{n} \\
  \cfalsify{\cand{\xi_1}{\xi_2}} &= \cand{\cfalsify{\xi_1}}{\cfalsify{\xi_2}} \\
  \cfalsify{\cor{\xi_1}{\xi_2}} &= \cor{\cfalsify{\xi_1}}{\cfalsify{\xi_2}} \\
  \cfalsify{\cinl{\xi}} &= \cinl{\cfalsify{\xi}} \\
  \cfalsify{\cinr{\xi}} &= \cinr{\cfalsify{\xi}} \\
  \cfalsify{\cpair{\xi_1}{\xi_2}} &= \cpair{\cfalsify{\xi_1}}{\cfalsify{\xi_2}}
\end{align*}
  \caption{Truify and Falsify Constraints}
  \label{fig:truify-falsify}
\end{figure}

\figurename~\ref{fig:truify-falsify} defines function $truify$ and $falsify$,
which substitute unknown constraint with truth constraint $\ctruth$ and falsity
constraint $\cfalsify$ respectively.

\begin{lemma}
  Assume $\isVal{e}$. Then $\csatisfyormay{e}{\xi}$ iff $\csatisfy{e}{\ctruify{\xi}}$.
\end{lemma}

\begin{lemma}
  $\csatisfy{e}{\xi}$ iff $\csatisfy{e}{\cfalsify{\xi}}$.
\end{lemma}

\begin{lemma}
  Assume $\isVal{e}$ and $\ctruify{\xi}=\xi$. Then $\cnotsatisfy{e}{\xi}$ iff $\csatisfy{e}{\cdual{\xi}}$.
\end{lemma}

\begin{theorem}
  $\csatisfy{\xi_r}{\xi_{rs}}$ iff $\csatisfy{\ctruth}{\cor{\cdual{\ctruify{\xi_r}}}{\cfalsify{\xi_{rs}}}}$.
\end{theorem}

\begin{lemma}
  $\csatisfyormay{e}{\xi}$ iff $\csatisfyormay{e}{\ctruify{\xi}}$.
\end{lemma}

\begin{lemma}
  Assume $\ctruify{\xi}=\xi$. Then $\csatisfyormay{\ctruth}{\xi}$ iff $\csatisfy{\ctruth}{\xi}$.
\end{lemma}

\begin{theorem}
  $\csatisfyormay{\ctruth}{\xi}$ iff $\csatisfy{\ctruth}{\ctruify{\xi}}$.
\end{theorem}

Now we use \textit{incon} judgment to further determine $\csatisfy{\ctruth}{\xi}$.

% !TEX root= pattern-paper.tex

\begin{figure}[bp]
\judgbox{\cincon{\Xi}}{}

\begin{mathpar}
\Infer{\CINCTruth}{
  \cincon{\Xi}
}{
  \cincon{\Xi, \ctruth}
}

\Infer{\CINCFalsity}{ }{
  \cincon{\Xi, \cfalsity}
}

\Infer{\CINCNum}{
  n_1 \neq n_2
}{
  \cincon{\Xi, \cnum{n_1}, \cnum{n_2}}
}

\Infer{\CINCNotNum}{ }{
  \cincon{\Xi, \cnum{n}, \cnotnum{n}}
}

\Infer{\CINCAnd}{
  \cincon{\Xi, \xi_1, \xi_2}
}{
  \cincon{\Xi, \cand{\xi_1}{\xi_2}}
}

\Infer{\CINCOr}{
  \cincon{\Xi, \xi_1} \\
  \cincon{\Xi, \xi_2}
}{
  \cincon{\Xi, \cor{\xi_1}{\xi_2}}
}

\Infer{\CINCInj}{ }{
  \cincon{\Xi, \cinl{\xi_1}, \cinr{\xi_2}}
}

\Infer{\CINCInl}{
  \cincon{\setof{\xi' | \cinl{\xi'} \in \Xi},\xi}
}{
  \cincon{\Xi, \cinl{\xi}}
}

\Infer{\CINCInr}{
  \cincon{\setof{\xi' | \cinr{\xi'} \in \Xi},\xi}
}{
  \cincon{\Xi, \cinr{\xi}}
}

\Infer{\CINCPairL}{
    \cincon{\setof{\xi_1' | \cpair{\xi_1'}{\xi_2'} \in \Xi},\xi_1}
}{
    \cincon{\Xi, \cpair{\xi_1}{\xi_2}}
}

\Infer{\CINCPairR}{
    \cincon{\setof{\xi_2' | \cpair{\xi_1'}{\xi_2'} \in \Xi},\xi_2}
}{
    \cincon{\Xi, \cpair{\xi_1}{\xi_2}}
}
\end{mathpar}

  \caption{Inconsistency of Constraints}
  \label{fig:incon}
\end{figure}

\begin{theorem}
  Assume $\ctruify{\xi}=\xi$. It is decidable whether $\cincon{\xi}$.
\end{theorem}

\begin{theorem}
  Assume $\ctruify{\xi}=\xi$. Then $\cincon{\cdual{\xi}}$ iff $\csatisfy{\ctruth}{\xi}$.
\end{theorem}
\section{Conclusion}
We can adapt our core calculus to fit in a more expressive setting.
\clearpage

\bibliographystyle{splncs04}
\bibliography{references.bib}
\end{document}
